\section{Limitations and Future Work}
\label{sec:limitations_future_work}

\subsection{Limitations}
The main limitations of this study are \textit{topology scale}, \textit{testbed fidelity}, \textit{replay-capacity} coverage, and \textit{forecasting scope}. While these constraints enable rigorous controlled experiments, they bound the generality of results and motivate validation across larger, more heterogeneous quantum-network testbeds. \textit{First}, all experiments use a controlled 4-node, 4-path internal topology with mid-range horizons, which enables clean attribution of algorithm, allocator, and replay-capacity effects but may not capture behavior in larger or more heterogeneous quantum-network testbeds. \textit{Second}, the current software-defined testbeds abstract away control-plane latency, node-level scheduling delays, memory contention, and timing-dependent decoherence effects, so deployment-level conclusions require further validation across higher-fidelity and more diverse testbeds. \textit{Third}, replay capacity is evaluated using two anchoring rules (\texttt{Tb}, \texttt{T}) and three scale factors ($s \in \{1,1.5,2\}$), which exposes the \texttt{capacity paradox} and motivates broader tail-risk characterization at larger horizons and topologies. \textit{Lastly}, predictive context currently uses ARIMA-based forecasting; more expressive temporal models (\eg RNNs and Transformers) with calibrated uncertainty quantification remain future work.
% \noindent\textit{Topology and scale.}
% All experiments are conducted on a controlled 4-node, 4-path topology with mid-range horizons. While this enables clean attribution of algorithm, allocator, and replay-capacity effects, larger and more heterogeneous quantum networks may alter both absolute efficiency and allocator--capacity interactions.

% \noindent\textit{Simulation fidelity and hardware constraints.}
% The current evaluation framework is simulation-driven and does not yet incorporate hardware-grounded constraints such as control-plane latency, scheduling delays, memory contention, or timing-dependent resource effects. Deployment-level conclusions therefore require validation on realistic quantum-network infrastructure.

% \noindent\textit{Replay-capacity coverage.}
% Replay capacity is evaluated using two anchoring rules (\small\texttt{Tb}, \small\texttt{T}) and three scale factors ($s \in \{1,1.5,2\}$). While sufficient to expose the capacity paradox, broader sweeps across larger horizons and topologies are needed to fully characterize long-term tail-risk behavior.

% \noindent\textit{Forecasting scope.}
% Predictive context currently relies on ARIMA-based forecasting. More expressive temporal models (\eg RNNs or Transformers) and calibrated uncertainty mechanisms remain unexplored.

\subsection{Future Work}

Future work extends the limitations found along four directions: \textit{expanded testbed stress testing}, \textit{cross-testbed standardization}, \textit{larger routing and allocation spaces}, and \textit{automated threat-aware adaptation}. \textit{First}, we will evaluate pursuit--neural and other hybrid models on larger quantum-network testbeds under operational stressors---control-plane latency, node-level scheduling delays, memory contention, and timing-dependent decoherence---to stress-test robustness
claims in higher-fidelity testbed settings. \textit{Second}, we will conduct standardized cross-testbed benchmarking across sparse, compact, dense, and mid-scale topologies using consistent run configurations, determining whether pursuit--neural robustness and the capacity paradox generalize across network structures. \textit{Third}, we will extend the framework to larger routing graphs, expanded action spaces, and hierarchical allocation strategies, studying allocator--capacity interactions under increasingly combinatorial conditions. \textit{Lastly}, we will replace the current heuristic meta-policy with a learned threat detector and automated configuration selector that dynamically adapts allocator and replay-capacity settings under changing attack patterns.

% \noindent\textit{Hardware-grounded validation.}
% A primary next step is evaluating pursuit-based contextual bandits on realistic quantum-repeater testbeds under operational constraints such as latency, scheduling delays, and imperfect physical operations.

% \noindent\textit{Standardized cross-testbed benchmarking.}
% We are conducting follow-up evaluations across sparse, compact, dense, and mid-scale external testbeds using standardized run configurations to determine whether pursuit-neural robustness and the capacity paradox generalize across network structures.

% \noindent\textit{Scaling topology and allocation.}
% Future work will extend the framework to larger routing graphs, expanded action spaces, and hierarchical allocation strategies to study allocator--capacity interactions under increasingly combinatorial routing conditions.

% \noindent\textit{Automated threat-aware adaptation.}
% We also plan to replace the current heuristic meta-policy with a learned threat detector and automated configuration selector capable of dynamically adapting allocator and replay-capacity settings under changing attack conditions.